\chapter{内容流水线}

\section{原来的问题：太封闭}

在本章所述工具出现之前，往知识库里添加内容必须手写特定格式的 Markdown，
而且知识点字段必须填一个\emph{已经存在}的编号。
这是个死循环：要加内容得先有知识点表，而知识点表写死了十五个。
拿到一本真实手册想录进去，第一步就卡住——
不知道该往哪个知识点挂，也不知道能不能新建。

现在反过来做：\emph{先喂原始资料，让知识点从语料里长出来。}

\begin{code}
python3 -m tools.ingest raw/
python3 -m tools.ingest raw/ --apply
python3 -m tools.kb_audit
python3 -m tools.kb_audit --quarantine
python3 -m evalkit.factcheck_run
\end{code}

依次为：切片归类质检并暂存、确认后写入知识库、交叉校验查库内矛盾、
把冲突条目移进隔离区、对外核实并产出复核队列。
支持 \cd{.txt .md .pdf .docx .csv .tsv .xlsx} 七种格式。
PDF 按页读取，位置标签带页码，这是评委最可能要求当场翻阅的定位方式。

\section{三道关口，各查各的}

三个工具的分工见表 \ref{tab:gates}。

\begin{table}[htbp]\centering\small
\caption{三道关口的分工}\label{tab:gates}
\begin{tabular}{@{}p{3.2cm}p{4.6cm}p{4.4cm}@{}}
\toprule
工具 & 查什么 & 看不见什么 \\
\midrule
\cd{tools.ingest} & 摄入质量：乱码、目录页、无句读、重复 & 内容对不对 \\
\cd{tools.kb\_audit} & 切片\emph{之间}是否自相矛盾 & 与外部资料是否一致 \\
\cd{evalkit.factcheck\_run} & 与\emph{外部资料}是否一致 & 库内两条是否打架 \\
\bottomrule
\end{tabular}
\end{table}

第二道必须单独做。单条格式完全合规、出处也真实的两条切片，可以互相打架——
同一个报警代码在甲页写“急停”、在乙页写“过载”，两条都能过入库门禁，
外部核实也可能各自找到支持的来源。\emph{只有把它们放在一起比才能发现。}

这类矛盾的危害是隐蔽的：检索时两条都会被召回，
生成的断言无论采信哪一条都能通过审核（因为确实“有依据”），
于是\emph{幻觉率照样是零，而输出的内容有一半是错的}。

\section{交叉校验的一处逻辑倒置}

\begin{keynote}{踩坑：用相似度当冲突的前提，方向反了}
交叉校验第一版写的是“相似度够高才比对”，结果漏掉了最严重的一类矛盾：
两条切片都在讲 SRVO-005，一条说是机器人超程，另一条说是伺服放大器过载，
相似度却只有 $0.175$——\emph{恰恰因为它们说的完全不同}。

正确的判据是：讨论对象相同就该比，比的是它们给出的\emph{结论}是否一致。
相似度只用来区分冲突的类型——高相似度加数值不同是“数值被改”，
低相似度是“含义不同”，后者往往更严重。
\end{keynote}

\section{剔除靠隔离，不靠删除}

自动剔除听起来干净，但删掉的是别人辛苦收集的资料。
判据有误、阈值偏严、误伤，都会导致内容悄悄消失而无人知晓。
因此不合格的条目一律移进隔离区并写明原因，
只有显式加参数才真正丢弃。

冲突类问题还有一条特殊规则：\emph{冲突双方中已核实的留下，未核实的隔离；
双方核实状态相同时两条都留下，只报告}。

\begin{keynote}{踩坑：隔离把已核实的一方也带走了}
第一版没有这条规则，结果注入的错误数据把已人工核实的两条一起带下水——
\emph{错误数据把正确数据挤出知识库，这是最糟糕的一种失败}。
双方状态相同时机器无从判断谁对，此时该做的是叫人来裁，
而不是随便挑一条丢掉。
\end{keynote}

\section{摄入产出的状态是 sourced}

摄入产出的状态是 \cd{sourced}：出处是文件名加位置加文件指纹，可以回溯原文。
这与 \cd{verified} 是两回事，理由已在 3.3 节说明。
\cd{tools.ingest} 永远不写 \cd{verified=true}，有测试盯着。

\begin{keynote}{踩坑：加个字段就把整个系统打死了}
端到端验收时发现，摄入工具给切片加了 \cd{origin} 字段，
而数据结构不认这个字段，于是\emph{入库之后整个系统起不来}——
检索、诊断、生成全部无法初始化。单元测试没有覆盖到，
因为它们从不真正入库再读回。

知识库是多个工具共同写入的，字段会随工具演进而增加。
读取端必须对未知字段容错，否则任何一个工具加字段都会把整个系统打死。
\end{keynote}

\bibliographystyle{gbt7714-2005}
\bibliography{refs}
